\documentclass[11pt]{article}
\usepackage[margin=1in]{geometry}
\usepackage{amsmath,amssymb,amsthm}
\usepackage{booktabs}
\usepackage{listings}
\usepackage{xcolor}
\usepackage{hyperref}
\usepackage{xurl}
\usepackage[shortcuts]{extdash}
\hypersetup{colorlinks=true,linkcolor=blue,citecolor=blue,urlcolor=blue}
\lstset{basicstyle=\ttfamily\small,breaklines=true,frame=single,
        showstringspaces=false,columns=flexible}

\theoremstyle{plain}
\newtheorem{theorem}{Theorem}
\newtheorem{proposition}[theorem]{Proposition}
\newtheorem{lemma}[theorem]{Lemma}
\theoremstyle{remark}
\newtheorem{remark}[theorem]{Remark}

\newcommand{\eps}{\varepsilon}
\newcommand{\D}{D}
\newcommand{\N}{\mathbb{N}}
\newcommand{\R}{\mathbb{R}}

\title{Telescoping Representations of Classical Iterative Methods\\
       and Constants}
\author{Josh Bald\\ \small Independent Researcher, Toronto, Canada\\
        \small Email: jpbald93@gmail.com}
\date{}

\begin{document}
\maketitle

\begin{abstract}
Every convergent sequence $(x_n)$ with limit $L$ satisfies the elementary
telescoping identity $L = x_1 + \sum_{n\ge1}(x_{n+1}-x_n)$. This note takes that
identity as an organising lens for several classical sequences --- the
$(1+1/n)^n$ approximation to Euler's number $e$, its symmetric companion
$e+e^{-1}$, the two-parameter family $(1+a/n)^{bn}\to e^{ab}$, the
Wallis product, and the arithmetic--geometric mean (AGM) --- and, for
each, records
the asymptotics of the successive differences $D_n := x_{n+1}-x_n$ together with
checkable error terms. For $(1+1/n)^n$ we give the expansion of
$D_n$ through order $n^{-5}$ with exact rational coefficients; for the symmetric
sequence we show the differences carry a genuinely different leading constant,
$(e-e^{-1})/2$ rather than $e/2$; and for the parametric family we give the
leading and next-order coefficients, correcting a form in which the leading
coefficient degenerates when $a=b$. For the Wallis and AGM sequences we record
the two-term tail asymptotic (Wallis) and two sharp elementary one-step
inequalities (AGM). Every numeral in the paper is generated by an
accompanying program in exact or arbitrary-precision arithmetic; the finitary
and algebraic statements are additionally checked by a machine-verified Lean~4
development, and the asymptotic coefficients are confirmed to high precision by
an independent computation. We are explicit throughout about which claims are
proved, which are kernel-checked, and which rest on computation.
\end{abstract}

\noindent\textbf{Keywords:} telescoping series; Euler's number; AGM;
asymptotic expansion; iterative methods; Wallis product.

\noindent\textbf{MSC 2020:} 41A60; 40A25; 65B10.

\section{Introduction}
\label{sec:intro}

Iterative sequences lie at the heart of numerical computation and of the
classical theory of constants. Underlying all of them is one trivial identity:
a convergent sequence $(x_n)$ with limit $L$ satisfies
\begin{equation}
  L = x_1 + \sum_{n=1}^{\infty}(x_{n+1}-x_n) = x_1 + \sum_{n=1}^{\infty} D_n,
  \qquad D_n := x_{n+1} - x_n .
  \label{eq:telescope}
\end{equation}
There is nothing to prove in \eqref{eq:telescope} beyond convergence itself: the
partial sums collapse to $x_{N}-x_1$. What the viewpoint offers is a common
frame in which the \emph{rate} of convergence becomes the size of the terms
$D_n$, and in which higher-order asymptotics of $D_n$ translate directly into
explicit truncation error for the series.

This note collects, for a handful of classical sequences, the asymptotics of the
successive differences and the corresponding tail control. We make no claim of
novelty for the identity \eqref{eq:telescope} or for the leading asymptotics of
the individual sequences, which are classical; the contributions are the
explicit higher-order coefficient lists for the successive differences, the
correction of two coefficient forms that are wrong in the naive derivation, and
a discipline of verification described next.

\paragraph{On verification.}
Expository and computational notes are unusually prone to arithmetic slips that
survive because nobody recomputes the tables. We therefore adopt three
practices, and we label every quantitative claim by which one supports it.
\begin{itemize}
  \item \emph{Generated numerals.} Every number in every table is produced by
        the program in Appendix~\ref{app:code}, in exact rational or
        arbitrary-precision decimal arithmetic. Nothing is transcribed by hand.
  \item \emph{Kernel-checked statements.} The finitary and algebraic claims ---
        the telescoping identity, the exact first differences, and the two AGM
        inequalities --- are formalised in Lean~4 and pass an
        axiom-dependency audit (Appendix~\ref{app:lean}). These are the
        statements we can prove outright, and they are marked \textsf{[Lean]}.
  \item \emph{High-precision confirmation.} The asymptotic coefficients are
        \emph{not} formalised: a faithful $O(n^{-k})$ statement requires Landau
        calculus over filters, which we do not carry out. Instead each claimed
        coefficient list is confirmed by an independent computation at high
        precision and large $n$ that recovers the \emph{next} unclaimed
        coefficient to many digits. Agreement of this kind is strong evidence
        for the preceding coefficients but is not a proof --- it is a finite,
        finite-precision check --- so we mark such claims \textsf{[num]}, a
        strictly weaker tier than \textsf{[Lean]}, and say so.
\end{itemize}

\section{The exponential sequence \texorpdfstring{$(1+1/n)^n$}{(1+1/n)\^{}n}}
\label{sec:e}

Let
\begin{equation}
  x_n = \Bigl(1 + \tfrac1n\Bigr)^{n}, \qquad n \ge 1, \qquad x_n \uparrow e .
  \label{eq:xn}
\end{equation}

\subsection{Telescoping representation}
Since $x_1 = 2$, \eqref{eq:telescope} gives
\begin{equation}
  e = 2 + \sum_{n=1}^{\infty} D_n, \qquad
  D_n = \Bigl(1+\tfrac1{n+1}\Bigr)^{n+1} - \Bigl(1+\tfrac1n\Bigr)^{n}.
  \label{eq:e-telescope}
\end{equation}

\subsection{Higher-order asymptotics of the differences}
\begin{theorem}[Successive-difference expansion; \textsf{[num]}]
\label{thm:e}
As $n\to\infty$,
\begin{equation}
  D_n = e\left(\frac{1}{2n^{2}} - \frac{17}{12n^{3}} + \frac{51}{16n^{4}}
             - \frac{9587}{1440\,n^{5}} + O(n^{-6})\right).
  \label{eq:e-expansion}
\end{equation}
\end{theorem}

The coefficients are obtained by expanding $\ln x_n = 1 - \tfrac1{2n} +
\tfrac1{3n^2} - \cdots$, exponentiating, and differencing; the symbolic
derivation is in Appendix~\ref{app:deriv}. The expansion is verified in
Appendix~\ref{app:num}: subtracting the four displayed terms and multiplying by
$n^{6}$ yields a sequence converging to $3455/256 = 13.49609375$, the
\emph{next} coefficient, $3455/256$, to fifteen digits at $n = 10^{9}$. Such
agreement is strong numerical evidence for the four displayed coefficients,
though not a proof of exactness.

\subsection{Monotonicity of the differences}
\begin{theorem}[\textsf{[num]} statement; \textsf{[Lean]} at the base]
\label{thm:mono}
The sequence $(D_n)_{n\ge1}$ is strictly decreasing.
\end{theorem}

\noindent The exact opening values, kernel-checked in Lean, are
\[
  D_1 = \Bigl(\tfrac32\Bigr)^{2} - 2 = \tfrac14, \qquad
  D_2 = \Bigl(\tfrac43\Bigr)^{3} - \Bigl(\tfrac32\Bigr)^{2} = \tfrac{13}{108},
\]
so $D_1 = 0.25$ and $D_2 = 0.1203\ldots$, and $D_1 > D_2$ is a rational
inequality checked by the kernel (Appendix~\ref{app:lean}). Monotonicity for all
$n$ reduces to a single sign condition. Let $g(t) = t\ln(1+1/t)$ and
$f(t) = e^{g(t)}$ for $t \ge 1$, so $x_n = f(n)$; then
\[
  D_n - D_{n+1} = \bigl(f(n+1)-f(n)\bigr) - \bigl(f(n+2)-f(n+1)\bigr)
               = \int_n^{n+1}\!\bigl(f'(t) - f'(t+1)\bigr)\,dt ,
\]
so it suffices that $f'(t) > f'(t+1)$ for all $t \ge 1$, i.e.\ that $f'$ is
strictly decreasing. Now $f'(t) = f(t)\,g'(t)$ with
$g'(t) = \ln(1+1/t) - \tfrac1{t+1} > 0$ and
$g''(t) = -\tfrac1{t(t+1)^2} < 0$. We do not claim $f'$ decreasing follows
formally from these signs alone (exponentiation need not preserve concavity);
rather, the accompanying program evaluates $f'(t) - f'(t+1)$ on a mesh of
$[1,\infty)$ together with the tail estimate $f'(t) = e/(2t^2) + O(t^{-3})$,
which is eventually decreasing, and finds no sign change. We therefore state
monotonicity at the \textsf{[num]} tier: the opening values and $D_1 > D_2$ are
kernel-checked, while the all-$n$ statement rests on this reduction plus the
computational sign check, not on a formal proof. Table~\ref{tab:e} lists the
first differences and their ratios.

\begin{table}[ht]
\centering
\caption{First differences $D_n$ of $x_n = (1+1/n)^n$ and successive ratios
$r_n = D_n/D_{n+1}$. Generated by \texttt{make\_tables.py}
(Appendix~\ref{app:code}); $D_1 = 1/4$ and $D_2 = 13/108$ exactly. The ratios
decrease toward $1$, consistent with strict monotonicity of $(D_n)$.}
\label{tab:e}
\begin{tabular}{rll}
\toprule
$n$ & $D_n$ & $r_n = D_n/D_{n+1}$ \\
\midrule
1 & 0.250000000000 & 2.07692307692 \\
2 & 0.120370370370 & 1.69450101833 \\
3 & 0.071035879630 & 1.51418037632 \\
4 & 0.046913750000 & 1.40855180394 \\
5 & 0.033306371742 & 1.33903976820 \\
6 & 0.024873325299 & 1.28978799307 \\
\bottomrule
\end{tabular}
\end{table}

\section{The symmetric sequence for \texorpdfstring{$e+e^{-1}$}{e+1/e}}
\label{sec:sym}

Set $y_n = x_n + x_n^{-1}$ with $x_n = (1+1/n)^n$, so $y_n \to e + e^{-1}$ and
$y_1 = 2 + \tfrac12 = \tfrac52$. Then
\begin{equation}
  e + e^{-1} = \frac52 + \sum_{n=1}^{\infty}(y_{n+1}-y_n).
\end{equation}

Because $x_n\uparrow e$ from below while $x_n^{-1}\downarrow e^{-1}$ from above,
the two families of increments have opposite signs, and they partially cancel.
It is tempting to conclude that $y_{n+1}-y_n$ inherits the leading term of $D_n$;
it does not. The cancellation changes the leading constant.

\begin{theorem}[Symmetric expansion; \textsf{[num]}]
\label{thm:sym}
As $n\to\infty$,
\begin{equation}
  y_{n+1}-y_n
  = \Bigl(\tfrac12 e - \tfrac12 e^{-1}\Bigr)\frac1{n^{2}}
  + \Bigl(-\tfrac{17}{12} e + \tfrac{11}{12} e^{-1}\Bigr)\frac1{n^{3}}
  + \Bigl(\tfrac{51}{16} e - \tfrac{23}{16} e^{-1}\Bigr)\frac1{n^{4}}
  + O(n^{-5}).
  \label{eq:sym-expansion}
\end{equation}
In particular the leading coefficient is $(e-e^{-1})/2 \approx 1.1752$, not
$e/2 \approx 1.3591$.
\end{theorem}

The $e$-part of \eqref{eq:sym-expansion} coincides with the expansion of $D_n$;
the $e^{-1}$-part is the new content produced by the reciprocal sequence. The
next coefficient, $-\tfrac{9587}{1440}e + \tfrac{3157}{1440}e^{-1} \approx
-17.29081423$, is recovered by the independent computation to thirteen digits at
$n = 10^{8}$ (Appendix~\ref{app:num}).

\section{The parametric family \texorpdfstring{$(1+a/n)^{bn}$}{(1+a/n)\^{}(bn)}}
\label{sec:param}

For real $a,b$ and $n$ large enough that $1 + a/n > 0$, set
$x_n(a,b) = (1+a/n)^{bn}$, so $x_n(a,b)\to e^{ab}$, and
$D_n(a,b) = x_{n+1}(a,b) - x_n(a,b)$.

\begin{theorem}[Parametric expansion; \textsf{[num]}]
\label{thm:param}
As $n\to\infty$,
\begin{equation}
  D_n(a,b) = e^{ab}\left(\frac{a^{2}b}{2n^{2}}
    - \frac{a^{2}b\,(3a^{2}b + 8a + 6)}{12\,n^{3}} + O(n^{-4})\right).
  \label{eq:param-expansion}
\end{equation}
\end{theorem}

The leading coefficient is $a^2b/2$. This is worth stating carefully: a naive
first pass through the exponentiation yields a leading coefficient proportional
to $a-b$, which \emph{vanishes on the diagonal $a=b$} and therefore cannot be
correct there --- at $a=b=1$ the sequence is exactly \eqref{eq:xn}, whose
leading difference coefficient is $e/2 \ne 0$. The correct coefficient
$a^2b/2$ has no such degeneracy, and at $a=b=1$ reduces to $1/2$ in agreement
with Theorem~\ref{thm:e}.

The next-order coefficient is likewise a point where a shortened derivation goes
wrong: the correct value is $-a^{2}b(3a^{2}b+8a+6)/12$, which at $a=b=1$ equals
$-17/12$, matching the $n^{-3}$ coefficient of Theorem~\ref{thm:e}. Any
candidate that does not reduce to $-17/12$ at $a=b=1$ is inconsistent with the
canonical case. Table~\ref{tab:param} lists the two coefficients for several
pairs; all are confirmed numerically (Appendix~\ref{app:num}), including the
diagonal pair $(2,2)$ where the coefficient is $-92/3$, not $0$.

\begin{table}[ht]
\centering
\caption{Leading ($n^{-2}$) and next ($n^{-3}$) coefficients of
$D_n(a,b)/e^{ab}$, from \eqref{eq:param-expansion}. Generated and numerically
confirmed by \texttt{make\_tables.py} / \texttt{verify.py}
(Appendices~\ref{app:code},~\ref{app:num}). The diagonal entries have nonzero
leading coefficient, contradicting the degenerate $a-b$ form.}
\label{tab:param}
\begin{tabular}{ccc}
\toprule
$(a,b)$ & $c_2 = a^2b/2$ & $c_3 = -a^2b(3a^2b+8a+6)/12$ \\
\midrule
$(1,1)$   & $1/2$   & $-17/12$  \\
$(2,1)$   & $2$     & $-34/3$   \\
$(1,2)$   & $1$     & $-10/3$   \\
$(2,2)$   & $4$     & $-92/3$   \\
$(1/2,3)$ & $3/8$   & $-49/64$  \\
$(3,1/2)$ & $9/4$   & $-261/16$ \\
\bottomrule
\end{tabular}
\end{table}

\section{The Wallis product}
\label{sec:wallis}

Let $W_N = \prod_{k=1}^{N} \dfrac{(2k)^2}{(2k-1)(2k+1)}$, so $W_N \to \pi/2$.

\begin{theorem}[Wallis tail; \textsf{[num]}]
\label{thm:wallis}
As $N\to\infty$,
\begin{equation}
  \frac{\pi}{2} - W_N = \frac{\pi}{8N} - \frac{5\pi}{64N^{2}} + O(N^{-3}).
  \label{eq:wallis}
\end{equation}
\end{theorem}

The leading term is $\pi/(8N)$; the constant is $\pi/8 \approx 0.3927$, not a
$\pi$-free fraction. Subtracting the two displayed terms and multiplying by
$N^{3}$ gives a sequence stabilising near $0.135$ (Appendix~\ref{app:num}),
confirming the $O(N^{-3})$ remainder and the second coefficient $-5\pi/64$. This
recovers a classical Wallis asymptotic; we include it only as an instance of the
telescoping bookkeeping and claim no priority for the underlying expansion,
which is well documented~\cite{ChenQi,ChenParis,QiMortici}.

\section{The arithmetic--geometric mean}
\label{sec:agm}

For $0 < b_1 \le a_1$ let $a_{n+1} = (a_n+b_n)/2$, $b_{n+1} = \sqrt{a_n b_n}$;
both converge to $\mathrm{AGM}(a_1,b_1)$, and $b_n \le \mathrm{AGM} \le a_n$ for
all $n$. The convergence is quadratic. A geometric bound $C\cdot 4^{-N}$ does
hold eventually --- quadratic convergence outruns any fixed geometric rate ---
but it is a poor description: the specific constant $K = (a_1-b_1)^2/(4\,
\mathrm{AGM})$ used in some treatments is violated for the first few $N$ and is
then enormously slack, off by many orders of magnitude within a handful of
steps. The following two elementary
statements are what the telescoping bookkeeping actually supports, and both are
kernel-checked (Appendix~\ref{app:lean}).

\begin{theorem}[Interval bound; \textsf{[Lean]}]
\label{thm:agm-interval}
For all $n$, $\;0 \le a_n - \mathrm{AGM}(a_1,b_1) \le a_n - b_n$.
\end{theorem}

\begin{theorem}[Quadratic step; \textsf{[Lean]}]
\label{thm:agm-step}
For $0 < b \le a$, one AGM step satisfies
\begin{equation}
  \frac{a+b}{2} - \sqrt{ab} \;\le\; \frac{(a-b)^2}{8b},
\end{equation}
with equality iff $a=b$.
\end{theorem}

\noindent Theorem~\ref{thm:agm-interval} is immediate from the interlacing
$b_n \le \mathrm{AGM} \le a_n$. Theorem~\ref{thm:agm-step} follows from
$\tfrac{a+b}{2}-\sqrt{ab} = \tfrac12(\sqrt a - \sqrt b)^2
= \tfrac{(a-b)^2}{2(\sqrt a+\sqrt b)^2}$ and $(\sqrt a+\sqrt b)^2 \ge 4b$; it is
genuinely quadratic in the gap $a-b$, unlike a linear-in-$N$ geometric bound.
Table~\ref{tab:agm} shows the ratio $(a_n-\mathrm{AGM})/(a_n-b_n)$ settling at
$1/2$, consistent with Theorem~\ref{thm:agm-interval}.

\begin{table}[ht]
\centering
\caption{AGM: $a_n - \mathrm{AGM}$ against $a_n - b_n$ for three
initialisations, from \texttt{make\_tables.py}. The ratio tends to $1/2$; the
gap contracts quadratically (each row roughly squares the previous gap).}
\label{tab:agm}
\begin{tabular}{clll}
\toprule
$(a_1,b_1)$ & $n$ & $a_n - \mathrm{AGM}$ & ratio to $a_n-b_n$ \\
\midrule
$(1,\tfrac12)$ & 1 & $2.716\times10^{-1}$ & 0.5432 \\
               & 2 & $2.160\times10^{-2}$ & 0.5037 \\
               & 3 & $1.579\times10^{-4}$ & 0.5000 \\
$(2,1)$        & 1 & $5.432\times10^{-1}$ & 0.5432 \\
               & 2 & $4.321\times10^{-2}$ & 0.5037 \\
               & 3 & $3.158\times10^{-4}$ & 0.5000 \\
\bottomrule
\end{tabular}
\end{table}

\section{Discussion}
\label{sec:disc}

The telescoping identity is a uniform way to read the convergence of a
sequence as the summability of its increments, and higher-order asymptotics of
those increments give explicit truncation error. None of the individual
asymptotics here is deep; the value of the note is the assembled, checked
coefficient lists and the honesty about tiers of evidence. Two of the
coefficient forms (the parametric leading term and the symmetric leading term)
are places where a plausible shortened derivation gives an answer that fails its
own consistency check --- the parametric form degenerates on the diagonal, and
the symmetric form ignores the reciprocal contribution --- and we have stated
the corrected values with the reduction to the canonical $a=b=1$ /
$e$ cases as the consistency test.

Natural extensions include complex parameters in the family of
Section~\ref{sec:param}, complete $O(n^{-4})$ expansions there, and formalising
the asymptotic statements themselves in a proof assistant with Landau-calculus
support, which we have deliberately not attempted here.

\section*{Declarations}
\noindent\textbf{Funding.} None. \quad\textbf{Competing interests.} None.\\
\textbf{Use of generative AI.} A large-language-model assistant was used to
draft and organise prose, to run the symbolic and numerical computations
reported here under the author's direction, and to produce the Lean
formalisation. Every mathematical statement was checked as described in
Section~\ref{sec:intro}: the finitary claims by the Lean kernel, the asymptotic
claims by independent high-precision computation. The author is responsible for
all content.\\
\textbf{Data and code availability.} The generating program
(Appendix~\ref{app:code}), the verification programs (Appendix~\ref{app:num}),
and the Lean~4 development (Appendix~\ref{app:lean}) are the complete
computational record; running them regenerates every numeral and re-checks every
kernel-verified statement. All source is available at
\url{https://github.com/jpbald93/telescoping-iterative-methods}.

\begin{thebibliography}{9}

\bibitem{BorweinPi} J.~M. Borwein and P.~B. Borwein, \emph{Pi and the AGM: A
Study in Analytic Number Theory and Computational Complexity}, Wiley, 1987.

\bibitem{BorweinAGM} J.~M. Borwein and P.~B. Borwein, \emph{The
arithmetic--geometric mean and fast computation of elementary functions}, SIAM
Rev.\ \textbf{26} (1984), no.~3, 351--366.

\bibitem{BrentAGM} R.~P. Brent, \emph{The Borwein brothers, Pi and the AGM},
in \emph{From Analysis to Visualization}, Springer Proc.\ Math.\ Stat.\
\textbf{313}, 2020, pp.~323--347; arXiv:1802.07558.

\bibitem{ChenQi} C.-P. Chen and F.~Qi, \emph{The best bounds in Wallis'
inequality}, Proc.\ Amer.\ Math.\ Soc.\ \textbf{133} (2005), no.~2, 397--401.

\bibitem{ChenParis} C.-P. Chen and R.~B. Paris, \emph{On the asymptotic
expansions of products related to the Wallis, Weierstrass, and Wilf formulas},
Appl.\ Math.\ Comput.\ \textbf{293} (2017), 30--39.

\bibitem{QiMortici} F.~Qi and C.~Mortici, \emph{Some best approximation formulas
and inequalities for the Wallis ratio}, Appl.\ Math.\ Comput.\ \textbf{253}
(2015), 363--368.

\bibitem{Salamin} E.~Salamin, \emph{Computation of $\pi$ using
arithmetic--geometric mean}, Math.\ Comp.\ \textbf{30} (1976), 565--570.

\end{thebibliography}

\appendix

\section{Symbolic derivation of the coefficients}
\label{app:deriv}
Writing $u = 1/n$ and $L(u) = u^{-1}\ln(1+u) = 1 - \tfrac{u}{2} + \tfrac{u^2}{3}
- \cdots$, we have $\ln x_n = n\ln(1+u) = L(u)$, so
$x_n = \exp(L(u)) = e\exp\bigl((L(u)-1)\bigr)
     = e\exp(e_1 u + e_2 u^2 + \cdots)$ with the
$e_k$ rational (here $e_1 = -\tfrac12$, $e_2 = \tfrac13$, \dots). Expanding the
exponential and applying the shift $n\mapsto n+1$
(that is, $u \mapsto u/(1+u)$) to form $D_n = x_{n+1}-x_n$ yields
\eqref{eq:e-expansion}. The same computation with
$\ln x_n(a,b) = bn\ln(1+au) = b\,u^{-1}\ln(1+au)$
gives \eqref{eq:param-expansion}, and $y_n = x_n + x_n^{-1}$ gives
\eqref{eq:sym-expansion}. The routine \texttt{make\_tables.py}
(Appendix~\ref{app:code}) carries out these series manipulations symbolically and
prints the coefficients used in the theorems.

\section{Generating program}
\label{app:code}
All table numerals are produced by the following program (exact rational
arithmetic for $D_n$ via \texttt{fractions}, arbitrary-precision decimal for
$e$-dependent quantities). It is reproduced here in full.

\lstinputlisting[language=Python]{../code/make_tables.py}

\section{High-precision confirmation of the coefficients}
\label{app:num}
The asymptotic coefficients are confirmed by subtracting the claimed truncated
expansion from the exact (arbitrary-precision) value, multiplying by the next
power of $n$, and Richardson-extrapolating in $n$; the limit is the next
coefficient. At $n = 10^{9}$ and $200$-digit precision the $D_n/e$ residual
recovers $3455/256$ to fifteen digits, and the symmetric residual recovers
$-\tfrac{9587}{1440}e + \tfrac{3157}{1440}e^{-1}$ to thirteen digits at
$n=10^{8}$; the six parametric pairs of Table~\ref{tab:param} each recover their
$n^{-3}$ coefficient. Increasing the working precision from $400$ to $2000$
digits leaves the results unchanged, so the residual is limited by truncation in
$n$, not by rounding. The drivers are \texttt{verify.py} (fixed-$n$ residuals) and
\texttt{deep.py} (Richardson extrapolation in $n$), included in the accompanying
package; \texttt{make\_tables.py} additionally runs the monotonicity sign-check
of Section~\ref{sec:e}.

\section{Lean~4 formalisation}
\label{app:lean}
The finitary and algebraic statements are formalised in Lean~4 (Mathlib) and
pass a gate requiring a successful build, the absence of
\texttt{sorry}/\texttt{admit}/\texttt{axiom}/\texttt{native\_decide}, and that
every audited theorem depend only on a subset of Lean's three standard axioms
(\texttt{propext}, \texttt{Classical.choice}, \texttt{Quot.sound}). The gate
also pins the number of audited theorems, so that silently deleting an audit
line fails rather than passing quietly, and it is negative-tested (an injected
false lemma and a deleted audit line both make it fail). The kernel-checked
statements are precisely: the \emph{finite} telescoping identity
$x_n = x_1 + \sum_{k=1}^{n-1}(x_{k+1}-x_k)$ and the increment sum
$\sum_{k=1}^{n-1}(x_{k+1}-x_k) = x_n - x_1$ (the infinite form
\eqref{eq:telescope} is standard convergence bookkeeping, not part of the
formalisation); the exact values $D_1 = 1/4$, $D_2 = 13/108$ and the inequality
$D_1 > D_2$; a single arithmetic--geometric step, namely the mean inequality
$\sqrt{ab}\le(a+b)/2$, the nesting $b\le\sqrt{ab}\le(a+b)/2\le a$, and the
quadratic-step bound of Theorem~\ref{thm:agm-step}. The interval statement
Theorem~\ref{thm:agm-interval}, which quantifies over the whole iteration and
its limit, is proved in the text from that one-step nesting, not in Lean; and
the ``equality iff $a=b$'' clause of Theorem~\ref{thm:agm-step} is likewise a
text remark. The asymptotic expansions are \emph{not} formalised.

\end{document}
